\documentclass{ximera}

\title{Oefeningen: Standaardvorm en eenheden}
\author{Daan Lipkens}

\begin{document}
\begin{abstract}
    Oefeningen op het omzetten van eenheden, beduidende cijfers en standaardvorm.
\end{abstract}
\maketitle

\section*{Oefeningen}
Noteer in onderstaande oefeningen je eindantwoorden steeds in de wetenschappelijke notatie! Let goed op het correcte aantal beduidende cijfers.

\begin{exercise}
    Zet om naar de eenheid in het vet tussen haakjes. Vul je antwoord in het eerste vak in en de correcte macht van 10 in het tweede vak.
    \begin{enumerate}
    \item $8{,}3 \, \mathrm{mm} = \answer[tolerance=0.01]{8.3} \cdot 10^{\answer{-3}} \, \textbf{m}$
    \item $34 \, \mathrm{km} = \answer[tolerance=0.01]{3.4} \cdot 10^{\answer{4}} \, \textbf{m}$
    \item $9{,}0 \, \mathrm{nm} = \answer[tolerance=0.01]{9.0} \cdot 10^{\answer{-7}} \, \textbf{cm}$
    \item $20 \, \mathrm{dm} = \answer[tolerance=0.01]{2.0} \cdot 10^{\answer{3}} \, \textbf{mm}$
    \item $0{,}79 \, \mathrm{m} = \answer[tolerance=0.01]{7.9} \cdot 10^{\answer{-3}} \, \textbf{hm}$
    \item $25{,}4 \, \mathrm{Mm} = \answer[tolerance=0.01]{2.54} \cdot 10^{\answer{7}} \, \textbf{m}$
    \item $33{,}4 \, \mathrm{cm} = \answer[tolerance=0.01]{3.34} \cdot 10^{\answer{2}} \, \textbf{mm}$
    \item $0{,}5 \, \mathrm{\mu m} = \answer[tolerance=0.01]{5.0} \cdot 10^{\answer{2}} \, \textbf{nm}$
    \end{enumerate}
    \begin{solution}
        Gebruik het verschil in machten van 10 tussen de prefixen:
        \begin{itemize}
        \item (c) van nm ($10^{-9}$) naar cm ($10^{-2}$) is een stap van $10^{-7}$. Dus $9{,}0 \cdot 10^{-7} \, \mathrm{cm}$.
        \item (d) van dm ($10^{-1}$) naar mm ($10^{-3}$) is een stap van $10^{2}$. $20 \cdot 10^2 = 2{,}0 \cdot 10^3 \, \mathrm{mm}$.
        \item (f) van Mm ($10^6$) naar m ($10^0$) is een stap van $10^6$. $25{,}4 \cdot 10^6 = 2{,}54 \cdot 10^7 \, \mathrm{m}$.
        \end{itemize}
    \end{solution}
\end{exercise}

\begin{exercise}
    Zet om naar de eenheid in het vet tussen haakjes.
    \begin{enumerate}
    \item $3{,}0 \, \mathrm{cm}^{2} = \answer[tolerance=0.01]{3.0} \cdot 10^{\answer{-4}} \, \textbf{m}^{\textbf{2}}$
    \item $8{,}5 \, \mathrm{dm}^{2} = \answer[tolerance=0.01]{8.5} \cdot 10^{\answer{-2}} \, \textbf{m}^{\textbf{2}}$
    \item $12{,}3 \, \mathrm{m}^{3} = \answer[tolerance=0.01]{1.23} \cdot 10^{\answer{7}} \, \textbf{cm}^{\textbf{3}}$
    \item $0{,}5 \cdot 10^{3} \, \mathrm{mm}^{3} = \answer[tolerance=0.01]{5.0} \cdot 10^{\answer{-1}} \, \textbf{cm}^{\textbf{3}}$
    \item $0{,}02 \, \mathrm{L} = \answer[tolerance=0.01]{2.0} \cdot 10^{\answer{-5}} \, \textbf{m}^{\textbf{3}}$
    \item $101 \, \mathrm{mm}^{3} = \answer[tolerance=0.01]{1.01} \cdot 10^{\answer{-2}} \, \textbf{cL}$
    \item $0{,}05 \, \mathrm{hm}^{3} = \answer[tolerance=0.01]{5.0} \cdot 10^{\answer{31}} \, \textbf{nm}^{\textbf{3}}$
    \item $5{,}123 \, \mathrm{\mu m}^{2} = \answer[tolerance=0.001]{5.123} \cdot 10^{\answer{-18}} \, \textbf{km}^{\textbf{2}}$
    \end{enumerate}
    \begin{solution}
        Let op dat je de machten van 10 ook kwadrateert (bij oppervlakte) of tot de derde macht verheft (bij volume). 
        \begin{itemize}
        \item (c) $1 \, \mathrm{m} = 10^2 \, \mathrm{cm}$, dus $1 \, \mathrm{m}^3 = (10^2)^3 = 10^6 \, \mathrm{cm}^3$. $12{,}3 \cdot 10^6 = 1{,}23 \cdot 10^7 \, \mathrm{cm}^3$.
        \item (e) $1 \, \mathrm{L} = 1 \, \mathrm{dm}^3 = 10^{-3} \, \mathrm{m}^3$. $0{,}02 \cdot 10^{-3} = 2{,}0 \cdot 10^{-5} \, \mathrm{m}^3$.
        \item (g) Van hm ($10^2$) naar nm ($10^{-9}$) is een factor $10^{11}$. Verheven tot de 3e macht is dit $10^{33}$. $0{,}05 \cdot 10^{33} = 5{,}0 \cdot 10^{31} \, \mathrm{nm}^3$.
        \end{itemize}
    \end{solution}
\end{exercise}

\begin{exercise}
    Zet om naar de eenheid in het vet tussen haakjes. 
    \begin{enumerate}
    \item $5{,}0 \, \mathrm{m/s} = \answer[tolerance=0.1]{1.8} \cdot 10^{\answer{1}} \, \textbf{km/h}$
    \item $75{,}0 \, \mathrm{km/h} = \answer[tolerance=0.01]{2.08} \cdot 10^{\answer{1}} \, \textbf{m/s}$
    \item $33{,}5 \, \mathrm{g/cm}^{3} = \answer[tolerance=0.01]{3.35} \cdot 10^{\answer{4}} \, \textbf{kg/m}^{\textbf{3}}$
    \item $0{,}52 \, \mathrm{kg/m}^{3} = \answer[tolerance=0.01]{5.2} \cdot 10^{\answer{-4}} \, \textbf{g/cm}^{\textbf{3}}$
    \item $1{,}25 \cdot 10^{6} \, \mathrm{cm/min} = \answer[tolerance=0.01]{7.50} \cdot 10^{\answer{2}} \, \textbf{km/h}$
    \item $15{,}0 \, \mathrm{mm/s} = \answer[tolerance=0.01]{1.50} \cdot 10^{\answer{-4}} \, \textbf{hm/s}$
    \item $100 \, \mathrm{g/m}^{3} = \answer[tolerance=0.01]{1.00} \cdot 10^{\answer{-7}} \, \textbf{kg/cm}^{\textbf{3}}$
    \item $3{,}4 \, \mathrm{kg/cm}^{3} = \answer[tolerance=0.01]{3.4} \cdot 10^{\answer{6}} \, \textbf{kg/m}^{\textbf{3}}$
    \end{enumerate}
    \begin{solution}
        Zet telkens de teller en de noemer afzonderlijk om.
        \begin{itemize}
        \item (a) $5{,}0 \cdot 3{,}6 = 18 = 1{,}8 \cdot 10^1 \, \mathrm{km/h}$.
        \item (c) $\frac{33{,}5 \cdot 10^{-3} \, \mathrm{kg}}{10^{-6} \, \mathrm{m}^3} = 33{,}5 \cdot 10^3 = 3{,}35 \cdot 10^4 \, \mathrm{kg/m^3}$.
        \item (e) $\frac{1{,}25 \cdot 10^6 \cdot 10^{-5} \, \mathrm{km}}{(1/60) \, \mathrm{h}} = 12{,}5 \cdot 60 = 750 = 7{,}50 \cdot 10^2 \, \mathrm{km/h}$.
        \end{itemize}
    \end{solution}
\end{exercise}

\begin{exercise}
    Geef onderstaande tijdstippen weer in uren, minuten en seconden.
    \begin{enumerate}
    \item $2{,}50 \, \mathrm{h} = \answer{2} \, \mathrm{h} \quad \answer{30} \, \mathrm{min} \quad \answer{0} \, \mathrm{s}$
    \item $0{,}25 \, \mathrm{h} = \answer{0} \, \mathrm{h} \quad \answer{15} \, \mathrm{min} \quad \answer{0} \, \mathrm{s}$
    \item $1{,}34 \, \mathrm{h} = \answer{1} \, \mathrm{h} \quad \answer{20} \, \mathrm{min} \quad \answer{24} \, \mathrm{s}$
    \item $3{,}33 \, \mathrm{h} = \answer{3} \, \mathrm{h} \quad \answer{19} \, \mathrm{min} \quad \answer{48} \, \mathrm{s}$
    \item $0{,}13 \, \mathrm{h} = \answer{0} \, \mathrm{h} \quad \answer{7} \, \mathrm{min} \quad \answer{48} \, \mathrm{s}$
    \item $4{,}44 \, \mathrm{h} = \answer{4} \, \mathrm{h} \quad \answer{26} \, \mathrm{min} \quad \answer{24} \, \mathrm{s}$
    \end{enumerate}
    \begin{solution}
        Vermenigvuldig het decimaal gedeelte van het uur steeds met 60 om de minuten te krijgen. Het nieuwe decimaal gedeelte van de minuten vermenigvuldig je opnieuw met 60 voor de seconden.
        \textbf{Voorbeeld (c):} $1{,}34 \, \mathrm{h}$. Dat is $1 \, \mathrm{h}$. De rest is $0{,}34 \cdot 60 = 20{,}4 \, \mathrm{min}$. Dat is $20 \, \mathrm{min}$. De rest is $0{,}4 \cdot 60 = 24 \, \mathrm{s}$.
    \end{solution}
\end{exercise}

\begin{exercise}
    Geef onderstaande tijdstippen weer in uur (wetenschappelijke notatie).
    \begin{enumerate}
    \item 2 uur en 30 minuten = $\answer[tolerance=0.01]{2.5} \cdot 10^{\answer{0}} \, \mathrm{h}$
    \item 3 uur, 40 minuten en 30 seconden = $\answer[tolerance=0.01]{3.68} \cdot 10^{\answer{0}} \, \mathrm{h}$
    \item 10 uur, 53 minuten en 5 seconden = $\answer[tolerance=0.01]{1.09} \cdot 10^{\answer{1}} \, \mathrm{h}$
    \item 39 minuten en 43 seconden = $\answer[tolerance=0.01]{6.62} \cdot 10^{\answer{-1}} \, \mathrm{h}$
    \item 5 seconden = $\answer[tolerance=0.01]{1.39} \cdot 10^{\answer{-3}} \, \mathrm{h}$
    \item 16 minuten = $\answer[tolerance=0.01]{2.67} \cdot 10^{\answer{-1}} \, \mathrm{h}$
    \end{enumerate}
    \begin{solution}
        Deel het aantal minuten door 60 en het aantal seconden door 3600. Tel deze waarden bij de uren op en werk af in wetenschappelijke notatie.
        \textbf{Voorbeeld (b):} $3 + \frac{40}{60} + \frac{30}{3600} = 3{,}675 \, \mathrm{h} \rightarrow 3{,}68 \cdot 10^0 \, \mathrm{h}$.
    \end{solution}
\end{exercise}

\begin{exercise}
    In kernfysica maakt men geregeld gebruik van de \AA ngström (symbolisch: $1 \text{ \AA } = 1 \cdot 10^{-10} \, \mathrm{m}$) en de femtometer (symbolisch: $1 \, \mathrm{fm} = 1 \cdot 10^{-15} \, \mathrm{m}$) om afstanden uit te drukken. Vul aan in wetenschappelijke notatie:
    \begin{enumerate}
    \item $1 \text{ \AA } = \answer{1.0} \cdot 10^{\answer{5}} \, \mathrm{fm}$
    \item $1 \, \mathrm{fm} = \answer{1.0} \cdot 10^{\answer{-5}} \text{ \AA}$
    \end{enumerate}
    \begin{solution}
        De \AA ngström is $10^5$ keer groter dan een femtometer (verschil tussen $10^{-10}$ en $10^{-15}$).
    \end{solution}
\end{exercise}

\begin{exercise}
    Rangschik onderstaande waarden van klein naar groot: 
    \[ A = 11{,}6 \, \mathrm{mg}, \qquad B = 1021 \, \mathrm{\mu g}, \qquad C = 0{,}0000006 \, \mathrm{kg}, \qquad D = 0{,}31 \, \mathrm{mg}. \]
    
    1 (Kleinste waarde): \wordChoice{\choice{A}\choice{B}\choice{C}\choice[correct]{D}} \\
    2: \wordChoice{\choice{A}\choice{B}\choice[correct]{C}\choice{D}} \\
    3: \wordChoice{\choice{A}\choice[correct]{B}\choice{C}\choice{D}} \\
    4 (Grootste waarde): \wordChoice{\choice[correct]{A}\choice{B}\choice{C}\choice{D}}
    
    \begin{solution}
        Zet alles eerst om naar dezelfde eenheid (bijv. kilogram) in wetenschappelijke notatie om ze makkelijk te vergelijken:
        \begin{itemize}
        \item $D = 0{,}31 \, \mathrm{mg} = 3{,}1 \cdot 10^{-7} \, \mathrm{kg}$
        \item $C = 0{,}0000006 \, \mathrm{kg} = 6 \cdot 10^{-7} \, \mathrm{kg}$
        \item $B = 1021 \, \mathrm{\mu g} = 1{,}021 \cdot 10^{-6} \, \mathrm{kg}$
        \item $A = 11{,}6 \, \mathrm{mg} = 1{,}16 \cdot 10^{-5} \, \mathrm{kg}$
        \end{itemize}
    \end{solution}
\end{exercise}

\begin{exercise}
    Bij het maken van oefeningen kom je onderstaande uitdrukkingen uit. Reken de uitdrukking uit en noteer de bijhorende eenheid van de uitkomst. Houd rekening met de regels van de beduidende cijfers en noteer in wetenschappelijke notatie.
    \begin{enumerate}
    \item $\dfrac{32{,}3 \, \mathrm{N}}{9{,}81 \, \mathrm{N/kg}} = \answer[tolerance=0.01]{3.29} \cdot 10^{\answer{0}} \, \textbf{kg}$
    \item $\dfrac{25{,}0 \, \mathrm{m}^{2} \cdot 51{,}0 \, \mathrm{kg/m}^{3}}{3{,}0 \, \mathrm{kg}} = \answer[tolerance=0.1]{4.3} \cdot 10^{\answer{2}} \, \textbf{m}^{\textbf{-1}}$
    \item $\dfrac{-9{,}81 \, \mathrm{m/s}^{2}}{2} \cdot (3{,}0 \, \mathrm{s})^{2} + 5{,}0 \, \mathrm{m/s} \cdot 3{,}0 \, \mathrm{s} + 50 \, \mathrm{m} = \answer[tolerance=0.1]{2.1} \cdot 10^{\answer{1}} \, \textbf{m}$
    \end{enumerate}
    \begin{solution}
        \begin{itemize}
        \item (a) Het quotiënt heeft 3 B.C. (beide factoren hebben 3 B.C.). $32{,}3 / 9{,}81 = 3{,}29 \, \mathrm{kg}$.
        \item (b) Het resultaat moet 2 B.C. hebben vanwege $3{,}0 \, \mathrm{kg}$. $(25{,}0 \cdot 51{,}0) / 3{,}0 = 425 \rightarrow 4{,}3 \cdot 10^2 \, \mathrm{m^{-1}}$. De eenheid is $\frac{\mathrm{m}^2 \cdot \mathrm{kg/m}^3}{\mathrm{kg}} = \frac{1}{\mathrm{m}} = \mathrm{m}^{-1}$.
        \item (c) Reken in stappen uit: $(-44{,}145 \, \mathrm{m}) + 15 \, \mathrm{m} + 50 \, \mathrm{m} = 20{,}855 \, \mathrm{m}$. Na afronden op eenheden (minst nauwkeurige term) krijg je $21 \, \mathrm{m} = 2{,}1 \cdot 10^1 \, \mathrm{m}$.
        \end{itemize}
    \end{solution}
\end{exercise}

\begin{exercise}
    Hoeveel beduidende cijfers hebben onderstaande meetresultaten? Geef voor elk meetresultaat het aantal beduidende cijfers.
    \begin{enumerate}
    \item 76,20 cm $\rightarrow \answer{4}$
    \item 0,302 kg $\rightarrow \answer{3}$
    \item 1,02 mm $\rightarrow \answer{3}$
    \item 63,00 mm $\rightarrow \answer{4}$
    \item $25 \cdot 10^{5}$ m $\rightarrow \answer{2}$
    \item 0,000 12 g $\rightarrow \answer{2}$
    \end{enumerate}
    \begin{solution}
        Voorloopnullen tellen \textbf{nooit} mee als beduidend cijfer. Nullen achteraan na de komma tellen \textbf{wel} mee, want ze geven de nauwkeurigheid van het meettoestel aan.
    \end{solution}
\end{exercise}

\begin{exercise}
    Rond onderstaande getallen af tot op \textbf{twee beduidende cijfers} (in wetenschappelijke notatie).
    \begin{enumerate}
    \item 3,85 m $= \answer[tolerance=0.01]{3.9} \cdot 10^{\answer{0}} \, \mathrm{m}$
    \item 10,0 mg $= \answer[tolerance=0.01]{1.0} \cdot 10^{\answer{1}} \, \mathrm{mg}$
    \item $4{,}96 \cdot 10^{5}$ m $= \answer[tolerance=0.01]{5.0} \cdot 10^{\answer{5}} \, \mathrm{m}$
    \item $53 \cdot 10^{23}$ kg $= \answer[tolerance=0.01]{5.3} \cdot 10^{\answer{24}} \, \mathrm{kg}$
    \item 0,041 32 cm $= \answer[tolerance=0.01]{4.1} \cdot 10^{\answer{-2}} \, \mathrm{cm}$
    \item 0,002 000 0 mm $= \answer[tolerance=0.01]{2.0} \cdot 10^{\answer{-3}} \, \mathrm{mm}$
    \end{enumerate}
\end{exercise}

\begin{exercise}
    Schrijf het resultaat met het juiste aantal beduidende cijfers (in wetenschappelijke notatie, houd overal je basiseenheid van de opgave vast).
    \begin{enumerate}
    \item $5{,}22 \, \mathrm{cm} - 3{,}6 \, \mathrm{cm} = \answer[tolerance=0.01]{1.6} \cdot 10^{\answer{0}} \, \mathrm{cm}$
    \item $2{,}1 \, \mathrm{cm} \cdot 200 \, \mathrm{cm} = \answer[tolerance=0.1]{4.2} \cdot 10^{\answer{2}} \, \mathrm{cm}^{\textbf{2}}$
    \item $6{,}00 \, \mathrm{m} + 5{,}7 \, \mathrm{cm} = \answer[tolerance=0.01]{6.06} \cdot 10^{\answer{0}} \, \mathrm{m}$
    \item $51{,}2 \, \mathrm{mm} \cdot 4 = \answer[tolerance=0.01]{2.05} \cdot 10^{\answer{2}} \, \mathrm{mm}$
    \item $\dfrac{19{,}62 \, \mathrm{g}}{3{,}0 \, \mathrm{mm}^{3}} = \answer[tolerance=0.01]{6.5} \cdot 10^{\answer{0}} \, \mathrm{g/mm}^{\textbf{3}}$
    \item $\dfrac{28{,}3 \, \mathrm{g}}{2} = \answer[tolerance=0.01]{1.42} \cdot 10^{\answer{1}} \, \mathrm{g}$
    \item $12 \cdot 10^{2} \, \mathrm{cm} \cdot 3{,}52 \, \mathrm{cm} = \answer[tolerance=0.1]{4.2} \cdot 10^{\answer{3}} \, \mathrm{cm}^{\textbf{2}}$
    \item $10{,}7 \, \mathrm{cm} - 5{,}2 \, \mathrm{mm} = \answer[tolerance=0.01]{1.02} \cdot 10^{\answer{1}} \, \mathrm{cm}$
    \item $7{,}32 \, \mathrm{m} \cdot 5{,}45 \, \mathrm{m} \cdot 0{,}84 \, \mathrm{m} = \answer[tolerance=0.1]{3.4} \cdot 10^{\answer{1}} \, \mathrm{m}^{\textbf{3}}$
    \item $\dfrac{3{,}2 \cdot 10^{3} \, \mathrm{mm}}{5{,}1 \cdot 10^{-2} \, \mathrm{mm}} = \answer[tolerance=0.1]{6.3} \cdot 10^{\answer{4}}$
    \end{enumerate}
    \begin{solution}
        \begin{itemize}
        \item (c) Tel niet zomaar waarden op met verschillende eenheden! Maak er eerst $6{,}00 \, \mathrm{m} + 0{,}057 \, \mathrm{m}$ van. Dit is een optelling, dus het minste aantal plaatsen na de komma telt: $6{,}057 \rightarrow 6{,}06 \, \mathrm{m}$.
        \item (d) en (f) De getallen 4 en 2 zijn hier exacte getallen (bijvoorbeeld een factor in een formule) en beperken het aantal B.C. niet! De meetwaarde (met respectievelijk 3 B.C.) is bepalend.
        \item (h) Ook hier eerst dezelfde eenheid maken: $10{,}7 \, \mathrm{cm} - 0{,}52 \, \mathrm{cm} = 10{,}18 \, \mathrm{cm}$. Afronden op 1 plaats na de komma $\rightarrow 10{,}2 \, \mathrm{cm}$.
        \end{itemize}
    \end{solution}
\end{exercise}

\end{document}